\scanpage{48}
% L1
Sous la forme \unsure{précédente|précédent}, on voit que l'on a mieux que des invariant dans $M^+_i(K)$. En fait, on a des
\[
\underline{\gamma}^{i}(X)\in\mathrm{Ob}\, M^+_{i\,\mathrm{bipure}}(K),
\]
objets semi-simples, définis de la façon suivante ($K$ parfait, ce qu'on peut supposer) $\gamma^i(X)\subset\underset{\substack{U\text{ ouvert}\\\text{dense}}}{\underrightarrow{\lim}}H^i(\unsure{U|M},\mathbb Z(0))$ est le plus grand sous-motif \uline{bipure} de type $i$. \deleted{[\unsure{Alors|la} le quotient par $\gamma^i(X)$} \inserted{C'est un foncteur en $X$} [où on l'anneau des fonctions rat. sur $X$, comme alg. sur $K$], de façon évidente
\diagram{\begin{tikzcd}[ampersand replacement=\&,row sep=small]
X'\arrow[d]\\X
\end{tikzcd}}
[NB si $X'\mathrel{\unsure{-|\to}}X$ n'est pas dominant, $X$ \unsure{irrid|irréd|irréd.}, alors $\gamma^i(X)\mathrel{-}\gamma^i(X')$ est nul]. Pour un hom \uline{dominant}, $\gamma^i(X)\longrightarrow\sigma^{i}(X')$ est \uline{injectif}. En effet, on peut supposer (prenant des modèles convenables) que $X'\mathrel{-}X$ est propre, \inserted{donc surjectif}, $X$ et $X'$ lisses, alors $H^i(X)\mathrel{-}H^i(X')$ est \uline{injectif} comme il est bien connu \dots

% L19
\uline{Question} \deleted{\unsure{la relation|Conditions}} Propriétés de ``fidélité'' du foncteur précédent ?? On aimerait que tout \deleted{hom} système d'hom $\underline{\deleted{\unsure{\beta}}\gamma}_i(\deleted{\unsure{K}}R)\mathrel{\unsure{-|\to}}\underline{\gamma}_i(R')$
% archive sheet number: 47
